\section{Evaluation}\label{sec:evaluation}
\subsection{Cases and Semantic Checks}
We combine public inputs with models constructed for this study. From the
Train v1.0 repository\footnote{\url{https://github.com/FTSRG/trainbenchmark/tree/v1.0}}
\cite{szarnyas_etal_2018}, we take its railway metamodel and first stored batch
snapshot, a generated benchmark model rather than operational railway data.
After adapting the metamodel to the supported Ecore subset, we import the
unchanged snapshot: ten classes, 21 features, two enumerations, 754 objects and
3,378 observation rows. Our generated textual rendition elaborates exactly to
that Core pair, including identities, flags, row order and absence. Removing all
but one required sensor from a Route gives a multiplicity failure while the
other structural checks pass.

We construct the other positive and negative cases to isolate typing,
multiplicity, inheritance, opposites and containment. The nine scaling inputs
in Section~\ref{sec:timing} and two small authored examples form the eleven
aligned positives in Table~\ref{tab:correctness}; they are not eleven independent
metamodels. An additional public input, the EMF Compare metamodel from its
3.3.28 repository, is rejected by the bridge for unsupported operations before
Core checking. The input and runner index is \texttt{experiments/README.md};
public originals are acquired separately, and historical measurement logs are
kept outside the source repository. The artifact records revisions, adaptations
and licenses.

For RQ1, Section~\ref{sec:conformance} supplies the source correspondence and
subset restrictions; these cases illustrate consequences of the predicates.
Section~\ref{sec:guarantees}'s proofs address RQ2. A Lean proof client derives
nonconformance from a repeated link with an upper-one opposite, exercising
Proposition~\ref{prop:opposite} under our equal-count restriction.

\subsection{Comparison with EMF}
We use EMF Ecore/XMI 2.39.0 and Common 2.42.0 under Java~21. An isolated registry selects EcoreValidator and generic
EObjectValidator for dynamic instances; Diagnostician checks schema and instance
roots separately. Generated validators and OCL delegates are excluded.
Table~\ref{tab:correctness}(a) separates these calls from VL-MOF's combined check:
instance acceptance cannot override schema rejection. Panel (b) reports the loaded-model comparison of Section~\ref{sec:tool}.
It uses values when native \texttt{eIsSet} holds and an empty sequence otherwise,
not unconditional \texttt{eGet}; equal states can still receive different decisions.
These comparisons measure agreement under the stated configurations; they do
not establish either validator's fidelity to EMOF.

\begin{table}[tbp]
\centering\small
\caption{Validation decisions (a) and loaded-state correspondence (b).}
\label{tab:correctness}
\textbf{(a) Validation decisions}\\[3pt]
\begin{tabular}{p{49mm}ccc}
\hline
Case & EMF schema & EMF instance & VL-MOF \\
\hline
Aligned positives (11) & Accept & Accept & Accept \\
Missing required String & Accept & Reject & Reject \\
Repeated unique String & Accept & Accept & Reject \\
Train with explicit defaults & Accept & Accept & Accept \\
Train with original omissions & Accept & Accept & Accept \\
\hline
Lower 1 / upper 0 & Reject & Reject & Reject \\
Empty $0..0$ & Reject & Accept & Accept \\
Repeated containment & Reject & Accept & Accept \\
Two unpaired features & Reject & Accept & Accept \\
\hline
\end{tabular}

\medskip
\textbf{(b) Correspondence after loading}\\[3pt]
\begin{tabular}{>{\raggedright\arraybackslash}p{54mm}>{\raggedright\arraybackslash}p{60mm}}
\hline
Case(s) & State comparison \\
\hline
All cases in (a) except those below & Match \\
Train with explicit defaults & 12 scalar-presence differences \\
Repeated containment; two unpaired features & Comparison fails: duplicate collector rows from repeated visits \\
\hline
\end{tabular}
\end{table}

The lower block of Table~\ref{tab:correctness}(a) exposes Ecore schema
restrictions. In particular, an empty $0..0$ property is permitted by our
structural interpretation but rejected by Ecore at schema level.
The repeated unique String survives loading but is accepted by this generic
validator and rejected by VL-MOF; it does not establish the behavior of every
EMF editor or generated validator. Repeated traversal of a contained object
produces duplicate collector rows in the two unpaired cases. Their alignment is
conservatively unsuccessful even though their distinct URI/value comparisons
agree; they are not timed. The two-feature fixture shares a schema containing
an additional nonunique composite \emph{children} declaration. Its schema
rejection therefore does not isolate the two-feature instance behavior.

Our adaptation of the Train v1.0 metamodel removes one generator annotation and maps EJavaObject
attributes to EInt (\emph{id}, \emph{length}) and EBoolean (\emph{active}). It uses
$0..1$ for \emph{route}, \emph{target}, \emph{entry} and \emph{exit}, versus one in
the article's diagram; Figure~\ref{fig:train} shows the evaluated metamodel.
We also create a copy of the public snapshot with omitted enum defaults
written explicitly. In this default-materialized copy, native
\texttt{eIsSet} reports twelve explicit first-enumeration values unset while
bridge import retains \texttt{FAILURE}. Both validators accept, despite singleton
versus empty observations. This case is excluded from paired timing. The original-snapshot control
with adapted Ecore instead aligns and is accepted by
both: zero lower bounds permit omissions without inserting defaults. This result
is specific to the adapted metamodel; the unadapted schema is unsupported.

\subsection{Validation Time}\label{sec:timing}
We generate three workload families ourselves to vary size and shape while
holding each family's schema fixed. A \emph{containment star} has one parent
and many children with inherited scalar features and ordered value lists
(11/101/501 objects). A \emph{Train projection} repeats our five-object
route/switch example, with optional bounds (2/20/100 groups). A
\emph{containment chain} nests each child inside its predecessor
(11/101/501 objects, depths 10/100/500). These are synthetic tests of our checker,
not executions of the Train Benchmark's query workload; the public 754-object
snapshot is not timed. The artifact records structural counts; a chain of $N$ objects has $3N-2$ value occurrences.

Each case uses five independent process trials, ten warmups and thirty measured
repetitions, with both tool orders and alternate traversal directions. Builds
precede timing. We measure combined schema/instance validation on preloaded
inputs, excluding loading, serialization and diagnostics. The artifact also
records VL-MOF schema-only and read/decode phases, plus command-wall times;
EMF schema-only and load/decode phases were not retained. Instance-only validation
was not measured and cannot be inferred by subtracting schema time.

Table~\ref{tab:validation-times} reports the median of five process medians
in milliseconds; the artifact retains each process median, their minimum and
maximum, and all measured samples. Rows count feature observations.
\begin{table}[t]
\centering
\caption{Preloaded combined schema--instance validation time in ms: pooled
median [first, third quartile] from five independent processes, each with 10 warmups and 30
retained samples. Rows are complete feature-observation rows.}
\label{tab:validation-times}
\scriptsize
\setlength{\tabcolsep}{3pt}
\begin{tabular}{lrrll}
\hline
Family/shape & Objects & Rows & EMF & VL-MOF \\
\hline
Containment star & 11  & 42   & 0.654 [0.452, 0.950] & 0.677 [0.664, 0.712] \\
Containment star & 101 & 402  & 0.725 [0.533, 0.956] & 9.664 [9.477, 9.974] \\
Containment star & 501 & 2002 & 1.582 [1.152, 2.116] & 125.713 [123.854, 127.518] \\
Train projection & 10  & 26   & 0.768 [0.464, 1.121] & 2.231 [2.194, 2.287] \\
Train projection & 100 & 260  & 0.944 [0.756, 1.257] & 24.840 [24.353, 25.412] \\
Train projection & 500 & 1300 & 1.726 [1.270, 2.362] & 169.457 [158.682, 171.603] \\
Containment chain & 11  & 33   & 0.515 [0.381, 0.723] & 0.154 [0.152, 0.164] \\
Containment chain & 101 & 303  & 0.828 [0.600, 1.056] & 109.293 [107.933, 111.918] \\
Containment chain & 501 & 1503 & 2.037 [1.556, 2.834] & timeout \\
\hline
\end{tabular}
\end{table}

Each completed cell contains 150 samples per tool. For the 501-object chain,
all five EMF processes completed and supply the reported 150 samples. Each
VL-MOF benchmark process exceeded its 120-second process limit before emitting
its buffered sample array, so there are no retained VL-MOF samples for that
cell. This bounds the whole 40-iteration benchmark process; it does not show
that one validation takes 120 seconds.

The results show sensitivity to graph shape as well as size. At 101 objects,
VL-MOF's chain median is about eleven times its star median, although the chain
has fewer observation rows. EMF remains below 3 ms at the pooled median for
these inputs. These descriptive measurements compare the current implementations,
whose data structures and runtimes differ; they are neither an asymptotic result
nor evidence of a general speed ranking. Loading, parsing, normalization,
diagnostics, and process startup lie outside the measured window.


For RQ3, 40 of 45 paired timing trials complete; five VL-MOF processes time out
on the deepest chain. Raw samples and excluded preliminary runs remain in the artifact.
